% Transcription — fonds Alexandre Grothendieck, shelfmark 27, batch 1 (pages 1-20).
% Established from the facsimile archives/batches/27/batch-01.pdf (a mirror of
% https://grothendieck.umontpellier.fr/27.pdf), per the skill
% transcribe-grothendieck. The folder is thirty-six pages; this is the first
% of two batches (pages 21-36 are batch 2, done separately).
%
% Pages skipped: 7 and 16, blank cream sheets (only show-through); 20, the
% upside-down lower half of a typed leaf on Nil(a), nilpotent subalgebras and
% subalgebras equal to their own normaliser, with the blanks for the gothic
% letters never inked and nothing in his hand — a typed verso with nothing of
% his, skipped by the settled rule, although its subject is close to the pencil
% notes of page 19. Page 1 is a tan cover inscribed in pencil, top right,
% « Divers [pour] XII à XIV » (the numerals underlined); it takes no \page
% marker and its words become the section heading, with a \note.
%
% Pages summarised: 2, 3, 4, 5, 6 — five consecutive leaves of a typescript,
% paginated 6 to 10 by the typewriter, headed in pencil on the first of them
% « [Co] Exp XIII » (the first word doubtful), and to all appearances a leaf
% run of Exposé XIII of SGA 3 (Proposition 1.3, Corollaires 1.4-1.6, Remarque
% 1.7, on a group G acting on V, a subscheme W, the morphism psi : X -> V and
% « points W-réguliers »). Following folders 55 and 56 and the brief for this
% folder (the typed text is published), the typed body is not re-set: each
% page gets one French \note saying which statements it carries, and his ink
% interventions — substitutions, insertions, deletions, the renumbered
% conditions, the marginal (3) — are given in full in a list. The gothic
% letters (n, m_a), the psi and phi, and all the implication arrows were
% inked by hand into blanks the machine left; this is said once per page and
% not item by item. Overstrikes made by the typist (xxx) are not his and are
% not recorded; typed interlinear insertions that he tied in with an ink
% bracket are recorded as such. Small inked repairs of single typed letters
% (« Corollaire », « rédiduel », « Theorem ») are recorded collectively.
%
% Pages transcribed from his hand: 8-15, 17-19. They are fast drafting, in
% ink (black, then blue from page 11 on), page 19 mostly in pencil; the formulas
% and the numbered statements carry, much of the connective prose does not and
% is left \ill{}, following folder 29 batches 8-9 and folder 56.
%
% What the batch is about. Complements to Exposés XI-XIV of SGA 3.
% Typescript, pp. 2-6: Prop. 1.3 and its proof — for G, V, W geometrically
% irreducible with the inequality (2), a chain of implications (3) among
% conditions (i)-(vi) (n = m_a, dim X = dim V, psi generically quasi-finite /
% unramified / étale, psi unramified / étale / smooth at (e*,a), M_a and N
% coinciding near e), all equivalent when V is normal at a; Cor. 1.4 (phi
% smooth at (e,a) when n = m_a and V normal at a); Cor. 1.5, the version
% independent of a choice of a in W(k); Cor. 1.6, psi birational iff ... ,
% with (i)-(iv) characterising a in U(k); Remarque 1.7 defining W-regular
% points. His hand adds, among other things, « génériquement » and « enfin
% ... normal en a » to the conditions of 1.3, « régulier en e » on p. 4,
% « séparable sur k » in 1.5, « V normal, k alg. clos » in 1.6.
% Pages 8-10 (« Complément XII », ink on tan): plan of a §6 « Le tore central
% maximal » (6.1-6.3), a cancelled 6.4 (decomposition of a smooth algebraic
% group over a perfect field via its maximal central torus R and a finite
% subgroup Z of R), Th. 6.4 « (Rosenlicht) » on invertible functions on X × Y
% (phi(x,y) = phi_1(x) phi_2(y)), its reduction to X, Y affine of dimension 1,
% the corollary that a morphism G -> G_m sending e to 1 on a smooth connected
% group is a homomorphism, Hom_k(G, G_m) of finite type, and a Proposition:
% the universal épimorphisme G -> H onto a group of multiplicative type,
% compatibility with base change when k is perfect, H a torus when G is
% connected and smooth, with a warning for k imperfect and a cochain exact
% sequence C(e,R) -> C(H,R) -> Hom_gr(H°, R). Pages 11-12 (« XII », « XI »):
% Cor. 5.8 (subgroups of multiplicative type extend across codimension >= 2
% over S normal), Cor. 5.9 (a subtorus of the generic fibre extends iff the
% reduced neutral components at codimension-one points are of multiplicative
% type; margin « rédiger mieux »), a cancelled Cor. 5.10, and a Corollaire
% 4.? for XI: a compatible family H(n) of subgroups of multiplicative type
% over a local S comes from a unique subgroup H, with a Remarque on the
% primes dividing the torsion of the type. Pages 13-15: Pic of affine
% algebraic groups — Pic(G) finite over k algebraically closed, zero if G_réd
% is solvable, Pic(G) = Pic(G/B)/Im D(T) for G smooth connected, Pic(G) = 0
% for G semisimple iff simply connected; over an arbitrary k, n Pic(G) = 0
% for n the degree of a suitable finite extension k' (conditions a)-d)),
% Pic(G) = Pic(G_{k~}) under Hom_{k~-gr}(G, G_m) = 0, and 3°) for G unipotent.
% Page 17: a table of implications and non-implications among conditions
% (i)-(xiv) of some statement not on these leaves, with counterexamples
% SL(2,k) and « GP(2,k) » in characteristic 2, (xii) = (x) + (vii),
% (xiii) = (xii) + (v). Page 18: a lemma — a nilpotent Lie algebra h over an
% infinite field acting on V with every rho(x) non-injective kills a nonzero
% v — proved by the weight decomposition, the principle of prolongation of
% identities and Engel. Page 19 (pencil): root computations in a Lie algebra
% g, L + k g_alpha, [X_-alpha, sum lambda_beta X_-beta], and a boxed remark
% that L + k g_-beta is its own normaliser but contains no Cartan subalgebra
% (reading doubtful), with an ink chain (xiv) => (i) => (iii) => (v).
%
% Notation fixed once for the batch: the gothic letters of the typescript as
% \mathfrak{n}, \mathfrak{m}_a (Lie algebras of N and M_a — the page itself
% says « M_a et N coïncident »); on p. 4, (ii)-(iii), the two letters stand
% in the reverse order, \mathfrak{m} = \mathfrak{n}_a, and are kept so, with
% a note. The typed ø read as \phi, his ψ as \psi. On pages 8-19: G_m for the
% multiplicative group, « t.m. » kept as he abbreviates « type multiplicatif »,
% G_{réd}, G^0_{réd}; \mathbf{D}(T) for the bold D of page 13 (group of
% characters); \overline{H}_y, and on page 14 \bar k and \tilde k as written.
%
% Readings to check first: « pour » on the cover; « Co » on p. 2; the title
% line of §6 on p. 8 (« Relations avec ... »); the whole cancelled 6.4; the
% hypotheses of Th. 6.4 (the letter before (x,y) is overwritten); the
% bracket after « épimorphisme » on p. 10 and the proof sketch at its foot;
% the relation between H(n) and H(m) on p. 12; « de ... unipotent nul »
% (pp. 14-15); and on p. 19 « est son propre normalisateur ... mais ne
% contient que des ... algèbres de Cartan ».
%
% Pass: Opus 5.5 (claude-opus-5-5), 2026-09-24 — first pass, unchecked against the pages by a human.

\documentclass[11pt,a4paper]{article}
% Le préambule commun à toutes les transcriptions du fonds Grothendieck.
%
% Il tient en une idée : **l'incertitude se compose, elle ne se résout pas.**
% Une transcription qui lisse un mot douteux a détruit la seule chose pour
% laquelle on la faisait — un lecteur doit pouvoir savoir, à chaque endroit,
% ce qui était sur le feuillet et ce qui a été ajouté. Les six macros qui
% suivent sont donc l'appareil critique, et non de la décoration.
%
% Les mêmes commandes sont reconnues par `scripts/render.mjs`, qui produit la
% vue de lecture du site. Une macro ajoutée ici doit l'être aussi là-bas,
% faute de quoi le rendu échouera bruyamment — ce qui est voulu : un rendu
% silencieusement faux est pire qu'un rendu absent.

\ProvidesPackage{grothendieck}[2026/08/08 Transcriptions du fonds Grothendieck]

\usepackage{amsmath,amssymb,amsthm}
\usepackage{mathrsfs}
\usepackage{stmaryrd}
% Les diagrammes commutatifs sont partout dans ce fonds : c'est la langue
% même de la géométrie algébrique de Grothendieck, et une transcription qui
% les rendrait en prose ne transcrirait rien.
\usepackage{tikz-cd}
% Français par défaut — c'est la langue des feuillets — mais l'anglais est
% chargé aussi : la traduction et le résumé sont des documents anglais, et la
% typographie française (espaces avant les deux-points) y serait une faute.
% Ces documents basculent par \selectlanguage{english} en tête de corps.
\usepackage[english,french]{babel}
\usepackage{fontspec}
\usepackage{geometry}
\geometry{a4paper,margin=25mm,marginparwidth=28mm}
\usepackage{marginnote}
\usepackage{xcolor}
% ulem, not soul. `soul`'s \st and \ul are built on a character-by-character
% scan that XeLaTeX's font machinery breaks: the document fails with an
% undefined control sequence rather than degrading. ulem does the same two jobs
% and is Unicode-engine safe. `normalem` stops it hijacking \emph, which the
% transcriptions use for Grothendieck's own emphasis.
\usepackage[normalem]{ulem}
\usepackage{hyperref}
\hypersetup{colorlinks=true,linkcolor=black,urlcolor=black,pdfborder={0 0 0}}

% --- Métadonnées du lot -----------------------------------------------------
% Reprises telles quelles par le script de rendu, qui les affiche en tête de
% la vue de lecture. Elles fixent ce que le fichier prétend couvrir : sans
% elles, une transcription flotte sans point d'attache dans un fonds de seize
% mille feuillets.

\newcommand{\folder}[1]{\gdef\grothFolder{#1}}
\newcommand{\batch}[1]{\gdef\grothBatch{#1}}
\newcommand{\leaves}[2]{\gdef\grothFirst{#1}\gdef\grothLast{#2}}
\newcommand{\foldertitle}[1]{\gdef\grothFoldertitle{#1}}
\gdef\grothFolder{?}\gdef\grothBatch{?}\gdef\grothFirst{?}\gdef\grothLast{?}\gdef\grothFoldertitle{}

% --- Le résumé --------------------------------------------------------------
%
% Il ouvre la lecture modernisée, sous le titre « Résumé », et il est composé
% en petit. Ce n'est pas un ornement : c'est le seul passage du document qui
% ne soit pas des mathématiques, et un lecteur doit pouvoir le reconnaître
% avant de l'avoir lu — puis le sauter s'il connaît déjà le sujet, ou s'y
% arrêter s'il ne le connaît pas. Le filet qui le suit marque la reprise du
% corps.

\newenvironment{resume}
  {\small\setlength{\parskip}{0.45em}}
  {\par\medskip\hrule\medskip}

% --- Mots-clés ---------------------------------------------------------------
%
% En anglais, en clôture du résumé de la lecture modernisée : le vocabulaire
% moderne sous lequel chercher ce que les feuillets construisent. Le manifeste
% du site (`npm run manifest`) les extrait pour étiqueter la cote entière —
% cette ligne est la source unique des tags, il n'y a pas de fichier à côté.
\newcommand{\keywords}[1]{\par\smallskip\noindent
  {\footnotesize\textsc{Keywords} --- \textit{#1}}\par}

% --- L'appareil critique ----------------------------------------------------

% Le numéro de feuillet, en marge. C'est l'ancre qui permet de régler un
% désaccord sur une formule en regardant *un* feuillet plutôt qu'une chemise
% de six cents. Le site s'en sert aussi pour tourner les pages du fac-similé
% quand on fait défiler la transcription.
\newcommand{\leaf}[1]{%
  \marginnote{\footnotesize\color{gray!70}\textsf{#1}}%
  \label{leaf:#1}\ignorespaces}

% Illisible : on ne devine pas. Un mot inventé qui se lit comme les autres est
% le pire résultat possible.
\newcommand{\ill}{\textcolor{red!60!black}{[\,\ldots\,]}}

% Lecture douteuse : le mot est donné, et signalé comme incertain.
\newcommand{\uncertain}[1]{\uline{#1}}

% Ajout de l'éditeur — une lettre manquante, un mot sous-entendu.
\newcommand{\add}[1]{\textcolor{blue!50!black}{[#1]}}

% Biffé par Grothendieck lui-même. Ce qu'il a rayé fait partie du document :
% une rature dit souvent où la pensée a changé de direction.
\newcommand{\struck}[1]{\sout{#1}}

% Note du transcripteur, sur l'état matériel du feuillet.
\newcommand{\note}[1]{\par\smallskip{\small\color{gray!60!black}\textit{[#1]}}\par\smallskip}

% Note marginale de Grothendieck, distincte du corps du texte.
\newcommand{\marginal}[1]{\par\smallskip\noindent\hspace{1em}%
  {\small\color{gray!60!black}#1}\par\smallskip}

% --- Titre ------------------------------------------------------------------

\AtBeginDocument{%
  \begin{center}
    {\large\bfseries Fonds Alexandre Grothendieck --- cote n\textsuperscript{o}~\grothFolder}\\[2pt]
    {\itshape\grothFoldertitle}\\[4pt]
    {\small feuillets \grothFirst--\grothLast\ (lot \grothBatch)}\\[2pt]
    {\footnotesize\color{gray!60!black}Transcription établie d'après le fac-similé numérisé
     par l'Université de Montpellier.\\
     Les lectures incertaines sont soulignées, les illisibles notées [\,\ldots\,],
     les ajouts entre crochets.}
  \end{center}
  \vspace{1em}}

\endinput


\folder{27}
\batch{1}
\pages{1}{20}
\dating{[vers 1962-1970]}
\watermark{Édition de démonstration}
\foldertitle{SGA D [SGA 3]. Compléments XII à XVII : tapuscrit annoté (s.d.), notes manuscrites (s.d.)}

\begin{document}

\section*{Divers \uncertain{pour} XII à XIV}

\note{inscrit au crayon en haut à droite du feuillet de garde (page 1), qui ne
porte rien d'autre ; les deux chiffres romains sont soulignés. Le feuillet ne
reçoit pas de numéro de page}

\page{2}

\marginal{\uncertain{Co} Exp XIII}\note{au crayon, en haut à droite, « XIII »
souligné, à gauche du numéro tapé 6}

\note{feuillet de tapuscrit paginé 6, repris à la main. Le texte tapé est
publié : on en donne la teneur et ses interventions, sans le recomposer. Il
porte l'énoncé de la \textbf{Proposition 1.3} — avec les notations
précédentes, $G$, $V$, $W$ géométriquement irréductibles et l'inégalité (2)
valable, les conditions (i) $\mathfrak{n} = \mathfrak{m}_a$ et $N$ lisse sur
$k$ ; (ii) $\mathfrak{n} = \mathfrak{m}_a$ et $\dim X = \dim V$ (cf. lemme
1.2) ; (iii), (iii bis), (iii ter) $\mathfrak{n} = \mathfrak{m}_a$ et
$\psi : X \to V$ génériquement quasi-fini, resp. non ramifié, resp. étale ;
(iv), (iv bis) $\psi$ non ramifié, resp. étale, en $(e^{*}, a)$ ; (v) $\psi$
lisse en $(e^{*}, a)$, $N$ lisse sur $k$ et (2) une égalité ; (vi) $M_a$ et
$N$ coïncidant au voisinage de $e$ — la chaîne d'implications (3), et le
début de la démonstration par le lemme 1.1 et 1.2. Les lettres gothiques
$\mathfrak{n}$, $\mathfrak{m}_a$, les $\psi$ et toutes les flèches
d'implication sont portées à la main dans les blancs de la frappe}

\begin{itemize}
  \item première ligne : « Avec les notations précédentes, suppos\add{ons}
  \struck{encore} » — la désinence « ons » à l'encre sur la frappe ; l'ajout
  tapé en interligne « et $G, V, W$ géométriquement irréductibles. » est relié
  au texte par un trait d'encre et une virgule, et le « c » de
  « considérons » est surchargé en capitale ;
  \item dans (v), « égalité. » écrit à la main au début de la ligne suivante,
  suivi d'un mot manuscrit biffé, \struck{\ill{}} ;
  \item \marginal{(3)} en marge gauche de la chaîne d'implications, qu'il
  écrit à l'encre :
  \[
    (\mathrm{v}) \Leftrightarrow (\mathrm{iv\ bis}) \Rightarrow
    (\mathrm{iii\ ter}) \Rightarrow (\mathrm{i}) \Leftrightarrow
    (\mathrm{ii}) \Leftrightarrow (\mathrm{iii})
    \Leftrightarrow (\mathrm{iii\ bis}) \Leftrightarrow
    (\mathrm{iv}) \Leftrightarrow (\mathrm{vi}) ;
  \]
  « ${}\Leftrightarrow (\mathrm{iii\ bis}) \Leftrightarrow{}$ » est ajouté en
  interligne au-dessus, entre (iii) et (iv), avec une accolade ; un crochet
  enserre « (v) » ; le « bis » tapé de (iii ter) est surchargé en « ter » ;
  \item la suite : « si $V$ est \struck{lisse sur $k$ et $W$ géométriquement}
  \add{génériquement} réduit, on a (iii \emph{ter}) $\Leftrightarrow$ (i), si
  \struck{de plus} \add{enfin} $V$ est \struck{géométriquement} normal
  \add{en $a$} \struck{\ill{}} toutes les conditions envisagées sont
  équivalentes. » — « génériquement », « enfin » et « en a » de sa main en
  interligne, « ter » et « (i) » surchargés à l'encre ;
  \item dans la démonstration, « (i) $\Rightarrow$ (iv) $\Leftrightarrow$
  (vi) », le dernier chiffre corrigé à l'encre ; « (iv) implique
  \struck{(iii bis)} que $\psi$ est génériquement non ramifié » ; « donc »,
  « $\geqslant$ », « $\leqslant$ », « on a » et une virgule après « 1.2. »
  portés ou repassés à l'encre.
\end{itemize}

\page{3}

\note{feuillet paginé 7. Fin de la démonstration de la proposition 1.3 :
(iv bis) implique $\mathfrak{n} = \mathfrak{m}_a$ et équivaut à (v) ; si $V$
est génériquement réduit, (iii bis) $\Rightarrow$ (iii ter) ; si $V$ est
normal en $a$, (iv) $\Rightarrow$ (iv bis) par SGA I 9.5 (ii). Puis le
\textbf{Corollaire 1.4} ($\mathfrak{n} = \mathfrak{m}_a$, $V$ normal en $a$ :
$\phi$ est lisse en $(e, a)$, donc $M_a$ est lisse sur $k$ en $e$ et $\psi$
lisse en $(e^{*}, a)$) et le début de sa démonstration par la suite
d'inégalités $\dim V + \mathrm{rang}\, \mathfrak{n} \leqslant \dim(G \times W)
\leqslant \dim_{(e,a)} \phi^{-1}(a) + \dim V \leqslant \dim \mathfrak{m}_a +
\dim V$, la théorie de la dimension et un théorème de Chevalley. Mêmes lettres
et flèches portées à la main}

\begin{itemize}
  \item « Si $V$ est \struck{réduit (ou seulement} génériquement réduit\add{,}
  \struck{\ill{}} » — la virgule à l'encre ; l'interligne tapé « comme » est
  pris dans la phrase ;
  \item « et non seulement \add{non} ramifié » : l'interligne tapé « non »
  relié par un signe d'insertion à l'encre ; un mot tapé biffé à l'encre
  après « en $(e^{*}, a)$, » ;
  \item Corollaire 1.4 : après « $\mathfrak{n} = \mathfrak{m}_a$, » l'ajout
  tapé en interligne « $V$ normal en $a$. » est tiré dans la ligne par une
  accolade d'encre ; un signe biffé avant « donc » ;
  \item le « rang » repassé, le dernier « $V$ » de la suite d'inégalités et
  « provient » surchargés à l'encre ; retouches d'une lettre (« Corollaire »,
  « rédiduel », « générique ») qu'on ne relève pas une à une.
\end{itemize}

\page{4}

\marginal{$e$}\note{à l'encre, dans la marge gauche, en regard de la première
ligne}

\note{feuillet paginé 8. Fin de la démonstration de 1.4 (troisième inégalité
par l'isomorphisme $\phi^{-1}(a) \simeq M_a$, lissité de $\phi$ par un
théorème de EGA IV) ; remarque que la démonstration invoque des résultats
délicats mais n'est pas utilisée ensuite ; puis l'énoncé du
\textbf{Corollaire 1.5}, indépendant du choix de $a$ : conditions (i), (i bis),
(i ter) ($\psi$ génériquement non ramifié, étale, lisse avec $N$ lisse et (2)
une égalité), (ii) à (iv) (il existe $a \in W(k)$ tel que …), et
les implications (ii) $\Leftrightarrow$ (iii) $\Leftrightarrow$ (iv)
$\Rightarrow$ (i) $\Leftrightarrow$ (i bis) $\Leftrightarrow$ (i ter),
toutes équivalentes si $k$ est algébriquement clos. Lettres gothiques,
$\psi$, $\in$ et flèches portés à la main}

\begin{itemize}
  \item « c'est une égalité si et seulement si $M_a$ est \add{régulier en $e$
  i.e.} lisse sur $k$ en $e$ i.e. $\phi^{-1}(a)$ est lisse sur $k$ en
  $(e,a)$ » — « régulier en $e$ i.e. » de sa main en interligne, relié par
  une accolade ;
  \item « Comme $V$ est \struck{réduit} normal donc réduit en $e$ » ; les
  interlignes tapés « s'ensuit (par un » et « la » (« de \add{la} théorie des
  schémas ») sont reliés à l'encre ; « \struck{sur} en $(e,a)$ » ; « n'aurons »
  surchargé ;
  \item Corollaire 1.5 : « \struck{Considérons les} et supposons $V$
  génériquement \struck{réduit. Alors les …} \add{séparable sur $k$.}
  Considérons les conditions suivantes » — « séparable sur k. » de sa main
  en interligne ;
  \item (i) : un mot biffé avant $\psi$ ; « (i ter) » surchargé sur « (i
  bis) » ;
  \item dans (ii) et (iii), les lettres gothiques portées à la main se lisent
  « $\mathfrak{m} = \mathfrak{n}_a$ », dans l'ordre inverse de la page 2 ; on
  transcrit tel quel\note{sic ?}.
\end{itemize}

\page{5}

\note{feuillet paginé 9. Démonstration du corollaire 1.5 (réduction au cas
$k$ algébriquement clos, $\psi$ dominant, choix d'un point $x \in X(k)$ où
$\psi$ est non ramifié et d'image régulière) ; puis l'énoncé du
\textbf{Corollaire 1.6} : pour que $\psi$ soit \emph{birationnel}, il faut et il
suffit que les conditions de 1.5 soient satisfaites et qu'il existe un ouvert
non vide $U$ de $V$ tel que tout $a \in U(k)$ soit contenu dans un unique
conjugué de $W$ par un élément de $G(k)$ ; et, pour $a \in V(k)$,
l'équivalence de (i) un unique conjugué de $W$ contient $a$, (ii) l'ensemble
de ces conjugués est fini non vide, (iii) le schéma $\psi^{-1}(a)$ a un point
isolé. $\psi$, $\phi$, $\in$ et flèches portés à la main}

\begin{itemize}
  \item « On aura $x = \phi(g^{*}, a)$ » : le $\phi$ à l'encre ;
  \item Corollaire 1.6 : « Soient $G, V, W$ comme dans 1.3 » — le 3 surchargé
  à l'encre, un mot biffé ensuite — et, en interligne, « (\add{$V$ normal,
  $k$} alg. clos.) », les quatre premiers mots de sa main devant l'interligne
  tapé « alg. clos. » ;
  \item « un élé\add{ment} de $G(k)$ », la fin du mot à l'encre ;
  \item dans (iii), les guillemets « » autour de « des conjugués de $W$ qui
  contiennent $a$ » sont à l'encre.
\end{itemize}

\page{6}

\note{feuillet paginé 10. Suite du corollaire 1.6 : si de plus $W$ est réduit
et si $U$ est le plus grand ouvert de $V$ au-dessus duquel $\psi$ est un
isomorphisme, les conditions équivalent aussi à (iv) $a \in U(k)$ ; la
démonstration par le \emph{Main Theorem} de Zariski et la platitude de
$p : X \to G/N$. Puis la \textbf{Remarque 1.7} : sous 1.6, un point
satisfaisant (i) à (iii) sera appelé point \emph{$W$-régulier} de $V$ ; sans
$k$ algébriquement clos, $a$ est $W$-régulier si $\phi^{-1}(a)$ est non vide
et fini sur $k(a)$ ; sous les conditions plus générales de 1.5, on introduit le
degré de projection $d$ de $\psi$ et le plus grand ouvert $U$ au-dessus duquel
$\psi$ est étale et fini. $\psi$ et flèches portés à la main}

\begin{itemize}
  \item « Lorsque de plus \add{$W$ est} réduit » : l'interligne tapé « W est »
  tiré à l'encre ;
  \item deux traits de crayon dans la marge gauche, sous « (iv) $a \in U(k)$ » ;
  \item « \add{On} not\add{era} que \struck{\ill{}} l'hypothèse que $\psi$ est
  birationnel » ; « que $W$ \add{soi}t réduit ou non » ; « Th\add{eo}rem »
  surchargés à l'encre ;
  \item Remarque 1.7 : « point » et « $W$-régulier » soulignés à l'encre ;
  la phrase « Cette notion manifestement est stable par extension du corps de
  base $k \to k'$, avec $k'$ algébriquement clos. » est barrée d'un trait
  d'encre, que longe une ligne festonnée dont le sens n'est pas sûr ;
  \item « est \add{non vide et} fini sur $k(a)$\add{,} \struck{\ill{}} » —
  interligne tapé relié à l'encre, virgule à l'encre, un mot manuscrit biffé ;
  \item « étale » et « fini » soulignés à l'encre ; « Lorsque $V$ est normal
  (\add{et} $X$ réduit (p.ex. $W$ réduit)), alors un $a \in V$ » — « et » de
  sa main devant l'interligne tapé, les parenthèses à l'encre ;
  « ra\add{tion}nel » surchargé.
\end{itemize}

\section*{Complément XII}

\note{titre de sa main, souligné, en tête de la page 8}

\page{8}

\textbf{6.} \emph{Le tore central maximal}. \emph{Relations avec \ill{}
\ill{} \uncertain{groupe}}

\emph{Tores quotients d'un groupe algébrique}

\textbf{6.1.} \emph{Définition en général}\qquad Compatibilité avec extension de
la \struck{\ill{}} base, \uncertain{descente}, \ill{} …

\textbf{6.2.} \emph{Existence lorsqu'il existe un centr. \uncertain{réductif}}, etc …\note{les deux lignes de droite sont enfermées par un trait
d'accolade qui les rattache à 6.1 et 6.2}

\textbf{6.3.} \struck{N.B. \ill{} \ill{}, \uncertain{des} groupes algébriques}
\emph{\uncertain{Cas} d'un groupe algébrique : un \struck{\ill{}} maximal
central} [N.B. le groupe algébrique \uncertain{n'est} \ill{} \ill{}
\uncertain{affine} …].

\marginal{\uncertain{lisse} \struck{\ill{}} \uncertain{affine}}\note{encadré à
droite, au-dessus de 6.4, relié par un trait à la ligne de 6.4}

\struck{\textbf{6.4.} Soit $G$ un groupe algébrique \uncertain{lisse} sur un
corps \uncertain{parfait}, \struck{alors} soit $R$ son tore central maximal.
Alors il existe un sous-groupe fini $Z$ de $R$, tel que $G_Z$ soit
\struck{\ill{} \ill{} \ill{}} \ill{} \ill{} de $R_Z$ et d'un \ill{} \ill{}
$H$. De plus, $H$ est tel que \ill{}}\note{tout ce paragraphe est encadré et
annulé par quatre longues diagonales et un trait sous sa dernière ligne ;
l'indice de $G_Z$, $R_Z$ se lit mal (peut-être $G/Z$, $R/Z$)}

\textbf{Th. 6.4} (\uncertain{Rosenlicht})\note{nom entouré, au crayon, dans
l'interligne au-dessus du théorème}. $X$, $Y$ géom. intègres sur un corps $k$,
\uncertain{munis} de pts $a$, $b$ \uncertain{rationnels} dans $k$. Soit
\[
  \varphi : X \times Y \to \mathbf{G}_m \quad \text{tel que} \quad
  \varphi(x, y) = 1,
\]
\note{la lettre devant « $(x,y) = 1$ » est surchargée et l'argument ne se lit
pas sûrement}
alors $\varphi$ est de la forme
\[
  \varphi(x, y) = \varphi_1(x)\, \varphi_2(y)
\]
i.e. \struck{$\varphi_1(\varphi_2)$} $u(w)$ \ill{} \ill{} \ill{}
\uncertain{produisant}\ill{} \uncertain{morphismes} $X \to \mathbf{G}_m$
($Y \to \mathbf{G}_m$).

\page{9}

\struck{On peut supposer que \struck{\ill{} \ill{}} $X$, $Y$ affines.}

Equivalent : \uncertain{ceci} : \struck{\ill{}}

Si $u(x, b) = 1$, $u(a, y) = 1$ identiquement, alors $u(x, y) = 1$
identiquement. Sous cette forme, on \uncertain{peut supposer} \struck{\ill{}}
$k$ alg. clos. \struck{Alors \ill{} \ill{} \ill{}} \struck{N.B. \ill{} \ill{}
\ill{}} \struck{aussi \ill{}} \struck{\ill{} implique $X$, $Y$ \ill{}}
\struck{\ill{} \ill{} \ill{}} ; \ill{} \ill{} \uncertain{généralement}
\ill{} : \struck{$u(x,y)$} \ill{} \ill{} $X \times Y \to \mathbf{G}_m$
\ill{} de la forme …\note{passage très raturé : plusieurs lignes biffées,
un encadré, et une ligne en vague qui relie la phrase conservée à la ligne
encadrée}

On \uncertain{peut} supposer $X$, $Y$ affines, et une réduction immédiate nous
ramène au cas où $X$, $Y$ \add{sont} de dim 1.

\textbf{Corollaire.} Kif-kif, $G$ \uncertain{remplacé} par un \struck{\ill{}}
\uncertain{groupe} \ill{}

Soit $T$ un tore.\note{encadré à droite}

\struck{$X \times Y \subset \overline{X} \times \overline{Y}$} \ill{}

\textbf{Corollaire.} \uncertain{Tout} \uncertain{morphisme}
\struck{$O(\overline{X} \times \overline{Y})$} \struck{d'un} $u : G \to
\mathbf{G}_m$, \struck{qui} $G$ \uncertain{groupe} \emph{lisse connexe}, qui
envoie $e$ dans \uncertain{$e$}, est un \uncertain{homomorphisme} de groupes.
\struck{\ill{} \ill{}}\note{« lisse connexe » souligné au crayon ; la ligne
$X \times Y \subset \overline{X} \times \overline{Y}$ et le
$O(\overline{X} \times \overline{Y})$ sont barrés d'un long trait de crayon}

\struck{\textbf{Corollaire 3}} \struck{\ill{}} \struck{\textbf{Proposition}
\uncertain{?} $\mathrm{Hom}(G, \mathbf{G}_m)$ \ill{}}
\struck{\ill{}} $\mathrm{Hom}_k(G, \mathbf{G}_m)$ \struck{\ill{} \ill{} \ill{}
de type fini.} \struck{Déf} \ill{} \uncertain{groupe} de type fini.

\page{10}

\textbf{Proposition} \uncertain{?} Il existe un \struck{\ill{}}
[\uncertain{épimorphisme} \ill{}]\note{au-dessus, en interligne : « $G$
\struck{\ill{}} algébrique » et « \uncertain{hom} »}
\[
  G \longrightarrow H
\]
$H$ de t.m., tel que \struck{$H$} \struck{\ill{}} tout \uncertain{hom.} de $G$
dans un groupe de t.m. se factorise de façon unique par $H$.

\begin{tikzcd}[column sep=small, row sep=small]
R \arrow[r, hook] \arrow[d, hook] & Z \arrow[d, hook] \\
T \arrow[r, hook] & G \arrow[d] \\
& \Gamma
\end{tikzcd}

\note{diagramme de la marge gauche : $R \subset Z$, $T \subset G$, les
inclusions verticales, une flèche descendant de $G$ vers une lettre lue
$\Gamma$ ; à droite de $Z$ un griffonnage et une flèche descendante non
lus}

b) Lorsque $k$ est \uncertain{parfait}, la formation de $H$ est compatible avec
\struck{$H$} extension de la base.

c) Lorsque $G$ est connexe \add{et lisse}\note{« et lisse » au crayon, en
interligne}, $H$ est un tore.

\begin{tikzcd}
G \arrow[r, "\text{épim}"] & H \\
G^{\circ}_{\mathrm{réd}} \arrow[u, no head] \arrow[r, no head] & H^{\circ} \arrow[u]
\end{tikzcd}

\[
  G/G^{\circ}_{\mathrm{réd}} \to H/H^{\circ}
\]
\note{le trait entre $G$ et $G^{\circ}_{\mathrm{réd}}$ porte une marque
biffée}

\textbf{N.B.} Si $k$ est \uncertain{imparfait}, \add{même} si $G$
\ill{} \ill{} \ill{} \uncertain{est} connexe \uncertain{et lisse}, la formation
de $H$ n'est \uncertain{en général} pas compatible \ill{} …

\struck{\textbf{Corollaire}} \textbf{Proposition}. $G$ un groupe
\uncertain{connexe} \struck{\ill{}} sur $k$, $R$ un \struck{groupe de t.m.}
\add{tore}. Alors \uncertain{les} \uncertain{extensions} \ill{} de $G$ par $R$
qui \uncertain{splittent} \ill{} \struck{$H^{i}_{\mathrm{Hochschild}}(G, R)
\simeq$} \ill{} \ill{} \uncertain{principaux} homogènes
\uncertain{splittent}.\note{la ligne biffée porte, sous $H$, le mot
« Hochschild »}

\[
  0 \to C^{\cdot}(e, R) \to C^{\cdot}(H, R) \longrightarrow
  \mathrm{Hom}_{\mathrm{gr}}(H^{\cdot}, R) \to 0
\]
\[
  H^{i}(\mathrm{Hom}_{\mathrm{gr}}(H^{\cdot}, R))
\]
\note{un $E$ biffé devant $\mathrm{Hom}$ ; les deux lignes sont fermées à
droite par un crochet}

En effet, \add{d'ext. centrale} \ill{} \ill{} \uncertain{cocycles}
\ill{}, \ill{} \ill{} \ill{} d'une \uncertain{extension} de $H$ par $R$, donc
\uncertain{splitte}.

\section*{XII}

\note{« XII », souligné, dans la marge gauche en tête de la page 11}

\page{11}

\textbf{Cor. 5.8.} Supposons $S$ \uncertain{normal} \add{\uncertain{noeth.}}.
\struck{Soit} \uncertain{irréductible}.

Soit $U$ un ouvert de $S$, tel que $S - U$ de codim $\geqslant 2$. Alors tt
\struck{\ill{}} ss-groupe de t.m. de $G|U$ provient d'un unique ss-groupe de
t.m. de $G$.

\textbf{Cor. 5.9.} Supposons $S$ \uncertain{normal} \ill{} \ill{} \ill{} de
\ill{} point $y$, soit $\overline{H}_y$ un \struck{sous-groupe de t.m.}
\add{sous-tore} de $G_y$. Pour que $H_y$ \uncertain{soit} la fibre
\uncertain{générique} d'un \struck{sous-groupe de t.m.} \add{tore}
$\overline{H}$ de $G$, il faut et il suffit que pour tout $s \in S$ tel que
\struck{$\dim(\overline{H}_y)$} $\dim \mathcal{O}_{S,s} = 1$,
$(\overline{H}_y)^{\circ}_{s, \mathrm{réd}}$ soit \struck{un tore.} de t.m.

\marginal{rédiger mieux}\note{écrit de biais dans la marge gauche, en regard de
la fin du corollaire 5.9, qu'un trait vertical embrasse}

\struck{\textbf{Cor. 5.10.} Sous les conditions précédentes, supposons que
$\overline{H}_y$ soit \ill{} un \emph{tore} \ill{} \ill{} \ill{} \ill{}
\ill{} \ill{} \ill{} \ill{} \ill{} \ill{} \uncertain{résiduel}) \ill{}
\ill{} utilise un}\note{annulé par quatre diagonales ; une ligne biffée
et un interligne au milieu ne se lisent pas}

\emph{Alors} …

\emph{Lemme} ---

\note{en bas à droite : « TSVP »}

\page{12}

\emph{Lemme} utilisé au corollaire \uncertain{nouveau} ?

\marginal{!!}\note{dans la marge gauche, en regard de la première ligne}

XI \struck{\ill{}}, \emph{Corollaire} 4.\ill{}\note{« XI » souligné ; le chiffre
après « 4. » est surchargé} \struck{\ill{} \ill{}}

$G$ lisse affine sur $S$ local \ill{} \ill{} \uncertain{fini} $S$.
\struck{\ill{}} $H(n)$ famille de ss-groupes de $G$, de t.m., $H(n)$ \ill{}
$H(m)$\note{une relation entre $H(n)$ et $H(m)$, indicée, dont le signe ne se
lit pas sûrement}, on suppose qu'il existe un ss-groupe de t.m. $H_s$ de $G_s$
tel que $H(n)_s = H_s$ pour tt $n$. Alors $\exists$ ss-groupe $H$ \add{de t.m.}
unique de $G$ tel que $H$ \ill{} $= H(n)$, $\forall n$.\note{« de t.m. » est
entouré en interligne ; le signe entre $H$ et $H(n)$ porte peut-être un indice}

\emph{Remarque.} \struck{Je utilise} \uncertain{Envisager} \ill{} \ill{}
\ill{} \ill{}, \ill{} \ill{} dans l'un des diviseurs premiers \ill{} \ill{}
$P \neq \emptyset$ ($P$ \ill{} \ill{} \ill{} premiers) \uncertain{prouvera}
que la torsion du type $H_s$ \uncertain{est} première à $P$. On utilise pour
ceci la variante de la démonstration précédente, \uncertain{où} $M$
\uncertain{remplacé} \ill{} \struck{\ill{}} sous-groupe de t.m. de type
$=$ type de $H_s$ [et \ill{} \struck{\ill{}} utilisera \ill{} \ill{} de
\uncertain{densité} de \ill{} dém. \ill{} la \uncertain{fermeture} \ill{}
]\note{un trait vertical dans la marge gauche embrasse la remarque}

\page{13}

\textbf{Proposition}. Soit $G$ groupe algébrique affine sur corps alg. clos
$k$. Alors $\mathrm{Pic}(G)$ est fini. Si $G_{\mathrm{réd}}$ est
\struck{lisse et} \emph{résoluble}, \struck{$\mathrm{Pic}(G)$ \ill{}}
$\mathrm{Pic}(G)$ est nul. Si $G$ est lisse, \struck{connexe},
\struck{$G$ est \ill{} $\mathrm{Pic}(G)$ \ill{} \ill{} \ill{}}
\[
  \mathrm{Pic}(G) \simeq \mathrm{Pic}(G/B)/\mathrm{Im}\, \mathbf{D}(T)
\]
[$B$ Borel, $T \subset B$ tore maximal,
\[
  \mathbf{D}(T) \to \mathrm{Pic}(G/B), \qquad \mathbf{D}(T) \simeq
  \mathrm{Hom}(B, \mathbf{G}_m)
\]
\uncertain{par} \uncertain{caractéristiques}]\note{le $\simeq$ est porté sous
$\mathbf{D}(T)$, verticalement ; plusieurs soulignements au crayon}

\note{dans la marge gauche, en regard :
$R \subset B \subset G \to G/B$, et au-dessous « $R \cap T$ » avec un
$\breve{T}$ dont le signe n'est pas sûr}

Si $G$ est \struck{\ill{}} \add{semi-simple},
$\mathrm{Pic}(G) = 0$ ssi $G$ est simplement connexe.

\note{à gauche de cet énoncé, une ligne ondulée enferme le diagramme
suivant, qu'il met en regard de deux colonnes de quotients}

\begin{tikzcd}
\mathrm{Pic}(G/B) \arrow[r, leftrightarrow] \arrow[d, no head, "\wr"'] & \mathbf{D}(T/R) \arrow[d] \\
\mathrm{Pic}(G/B) & \mathbf{D}(T) \arrow[l]
\end{tikzcd}

\[
  \begin{array}{ll}
    G/R \to G/B, & T/R \subset G/R \\
    G \to G/B, & T \subset G
  \end{array}
\]
\note{colonne de gauche ; une flèche montante relie les deux lignes ; devant
$T/R$, une flèche partant d'un $T$ biffé}

\[
  G/R = GP(n) \to \text{Dém.}, \qquad
  G = GL(n) \to \text{Dém.}, \quad T(n) \subset G
\]
\note{« Dém. » est une lecture douteuse ; un $G$ biffé entre les deux lignes ; les deux exemples sont écrits l'un sous l'autre dans la marge gauche,
sous un trait}

\textbf{Proposition}. $G$ \struck{\ill{} \ill{}} groupe alg. affine sur un corps
$k$. Alors $\exists$ entier $n > 0$ [égal au degré d'une extension convenable
$k'$ de $k$] tel que \struck{\ill{}} $n \,\mathrm{Pic}(G) = 0$.
\struck{Si $G_{\mathrm{réd}}$ \ill{} \ill{}, alors \ill{} \ill{} \ill{}
pour un \ill{} \ill{} \ill{} du corps} \ill{} $G$. Les \uncertain{précisions}
\ill{} \ill{}\note{le passage biffé est encadré et barré de lignes en
festons}

\page{14}

1°) On peut prendre pour $n$ le degré de toute extension finie $k'$ de $k$ ayant
les propriétés suivantes :

a) $G_{k' \mathrm{réd}}$ est lisse sur $k'$ i.e. séparable$/k'$

b) \struck{$G_{k' \mathrm{réd}}$ admet un Borel} Les pts de $G/G^{\circ}$
sont \struck{définis} \add{rationnels} sur $k'$.

c) « Le » radical $R'$ \struck{est} \add{de $G^{\circ}_{k' \mathrm{réd}}$} est
défini sur $k'$, et splitté \struck{$G^{\circ}_{k' \mathrm{réd}}$ admet un
Borel $B'$ qui est splitté} sur $k'$ [suite de composition : quotients
successifs $\mathbf{G}_a$ ou $\mathbf{G}_m$]\note{« de $G^{\circ}_{k'
\mathrm{réd}}$ » est la ligne biffée d'abord écrite pour c), rattachée par un
trait au radical}

d) \add{\uncertain{hyp} \ill{}} $G^{\circ}_{k' \mathrm{réd}}/R'$ a
\struck{un groupe \ill{} \ill{}} \add{déployé} \struck{a un tore maximal de
$B'$} [\ill{} \uncertain{maximal} \uncertain{est} splitté] et $N(T')$,
$C(T')$ \struck{\ill{} \ill{} \ill{}! \ill{}, donc}\note{interlignes et
biffures enchevêtrés ; lecture d'ensemble douteuse}

\struck{Par suite, 1°) Si $G$ est unipotent \add{connexe}, on peut prendre \ill{} une
puissance de la caractéristique. 2°) Si $G$ est \struck{réductif} \ill{} simple, alors
($\mathrm{Hom}(G, \mathbf{G}_m)$ étant nul, après \ill{} \ill{}
\ill{}) \ill{} \ill{} \ill{}] $\mathrm{Pic}$ \struck{$G$ est lisse}}\note{le
bloc est encadré et annulé par trois diagonales}

2°) Supposons que $\mathrm{Hom}_{\bar k\text{-gr}}(G_{\bar k}, \mathbf{G}_{m
\bar k}) = 0$, et que $G$ est \add{lisse et} connexe\note{« lisse et » entouré,
en interligne, avec un « $k$-gr » au-dessus}. Soit $\tilde k$ la clôture
\add{algébrique} \uncertain{séparable} de $k$. Alors
$\mathrm{Pic}(G) \simeq \mathrm{Pic}(G_{\tilde k})$.
\struck{En \uncertain{particulier}, si $\mathrm{Hom}(G_{\bar k}, \mathbf{G}_m)$}
Si \uncertain{même} $\mathrm{Hom}_{\tilde k\text{-gr}}(G_{\tilde k},
\mathbf{G}_{m k}) = 0$ [ce \uncertain{qui} \ill{} \ill{} \ill{} de la
condition précédente \uncertain{sur} $\tilde k$, $G$ est de \ill{}
unipotent nul] … $= \mathrm{Pic}(G) \simeq \mathrm{Pic}(G_{\tilde k})$.
\note{« séparable » est souligné deux fois et « algébrique » écrit
au-dessus, sans que l'un biffe clairement l'autre ; un trait vertical dans la
marge embrasse le 2°)}

\page{15}

3°) Si $G$ est \emph{unipotent}, alors on \uncertain{sait} que
$\mathrm{Pic}(G)$ est annulé par une puissance de la car. $p$. Si $G$ est de
\ill{} unipotent nul et $\mathrm{Hom}_{\bar k\text{-gr}}(G_{\bar k},
\mathbf{G}_{m \bar k}) = 0$, alors $\mathrm{Pic}(G)$ \struck{fini}
$\simeq \mathrm{Pic}(G_{\bar k}) =$ \emph{groupe fini}.

\page{17}

\struck{(vii) $\overset{?}{\Longrightarrow}$ (vi)}

\struck{(vii) $\not\Rightarrow$ (vi) = si $\psi$ \emph{birationnel},
$\mathrm{Sl}(2, k)$ \emph{car. 2}.}

\struck{(vi) $\not\Rightarrow$ (v)}

\struck{(x) $\not\Rightarrow$ (i)}

\struck{(vii) $\not\Rightarrow$ (x) ok.}

\struck{a fortiori (vi) $\not\Rightarrow$ (i), (x) $\not\Rightarrow$ (xiii)}

\struck{(x) $\not\Rightarrow$ (i) ; Cas semi-simple --- --}

\struck{(i) $\not\Rightarrow$ (x) ; a fortiori (vi) $\not\Rightarrow$ (v),
(vi) $\not\Rightarrow$ (x), (i) $\Rightarrow$ (xiii) ; $GP(2, k)$,
car. 2}\note{toute la moitié supérieure de la page est annulée par de
longues diagonales ; un « ? » en marge gauche en regard de « (vi)
$\not\Rightarrow$ (v) » ; les groupes « a fortiori » sont entourés}

\marginal{!}\note{en marge gauche, avec un trait qui mène à l'accolade
suivante}
\[
  \left\{
  \begin{array}{ll}
    (\mathrm{vii}) \not\Rightarrow (\mathrm{vi}) & = \text{si } \psi
      \text{ birationnel} : \mathrm{Sl}(2, k), \text{ car. } 2 \\[3pt]
    (\mathrm{x}) \not\Rightarrow (\mathrm{v}) & [\mathrm{x} \not\Rightarrow
      (\mathrm{i}),\ (\mathrm{x}) \not\Rightarrow (\mathrm{xiii}),\
      (\mathrm{vi}) \Rightarrow (\mathrm{v})] \\
      & \text{Cas d'un système de } G \text{ semi-simple} \ldots \\[3pt]
    (\mathrm{i}) \not\Rightarrow (\mathrm{x}) & [(\mathrm{v}) \not\Rightarrow
      (\mathrm{x}),\ (\mathrm{vi}) \not\Rightarrow (\mathrm{x}),\
      (\mathrm{i}) \not\Rightarrow (\mathrm{xiii})] \\
      & GP(2, k), \text{ car. } 2
  \end{array}
  \right.
\]
\note{« système » est une lecture douteuse}

\note{dans la marge gauche, en regard, « $G = PGl(2, k)$ », « $\dim W = 1$ »,
« $\dim X : 3 - 1 + 2 = 4$ », un zigzag de biffure sur ces lignes, et une
boucle qui enferme}
\[
  (\mathrm{vii}) \not\Rightarrow (\mathrm{vi}), \qquad
  (\mathrm{x}) \not\Rightarrow (\mathrm{v}), \qquad
  (\mathrm{v}) \not\Rightarrow (\mathrm{i})
\]

(v) $\not\Rightarrow$ (i) ???\note{encadré} ; à droite, $\mathrm{x} \neq \mathrm{i}$

$\mathfrak{h} = \mathfrak{n}$\note{entouré}

(xii) $=$ (x) $+$ (vii)

(xiii) $=$ (xii) $+$ (v)\note{encadré ; le dernier signe, sous une tache
d'encre, ne se lit pas}

\page{18}

\struck{Si $\mathfrak{h}$ \ill{}}

\emph{Lemme}. $\mathfrak{h}$ alg. de Lie nilpotente \add{sur un corps
\emph{infini} $k$} opérant sur un vectoriel $V$. Supposons que pour tt
$x \in \mathfrak{h}$, $\rho(x)$ soit non injectif. Alors $\exists\, v \in V$,
$v \neq 0$ tel que $\rho(\mathfrak{h}) v = 0$.\note{« sur un corps infini $k$ »
est entouré, en interligne, relié au texte ; un trait vertical dans la marge
embrasse l'énoncé}

\uncertain{$\bar k$} \ill{}, $V = \coprod V_i$, $V_i$ \struck{propre}
\add{primaires} pour la fonction poids $\lambda_i$ sur $\mathfrak{h}$.
L'hypothèse signifie que pour tt $x \in \mathfrak{h}$, $\exists\, i$ tel que
$\lambda_i(x) = 0$, donc $V = \bigcup V(\lambda_i)$. Par le principe du
prolongement des identités, on en conclut que $\exists\, i$ tel que
$\lambda_i = 0$. Donc, $\mathfrak{h}$ opère sur $V_i$ \struck{\ill{}} par
opérations nilpotentes. Donc \struck{\ill{} \ill{}} par Engel,
$\exists\, v \in$ \uncertain{$\mathfrak{h}$}\note{sic ?}, $v \neq 0$, tel que
$\mathfrak{h} . v = 0$, ok.

\[
  P_i(x, t) = \sum c_i(x)\, t^{n-i}
\]
\[
  \underbrace{t^{n} - c_1(x)\, t^{n-1} \ldots \pm c_n}_{[t - \lambda(x)]^{n}} = 0
  \qquad \lambda(x)^{n} = c_n
\]
\note{ces trois lignes sont au crayon ; la lettre de $P_i$ et le signe du
dernier terme sont douteux}

\page{19}

\note{page au crayon, sauf la chaîne d'implications, à l'encre}

\[
  \mathfrak{L} \subset \mathfrak{g}, \qquad \mathfrak{L} + k\,
  \mathfrak{g}_{\alpha}, \qquad \sum_{\beta \neq \alpha} \lambda_{\beta}
  X_{-\beta}
\]
\struck{\ill{}} \quad $\bar\beta = 0$
\[
  \Bigl[ X_{-\alpha}, \sum_{\beta \neq \alpha} \lambda_{\beta} X_{-\beta}
  \Bigr] = \sum \lambda_{\beta} [X_{-\alpha}, X_{-\beta}], \qquad
  [X_{\alpha}, X_{\beta}]
\]

Deux racines \qquad $X_1, X_2, X_3$

\note{suivent deux esquisses de matrices, une $2 \times 2$ et une $3 \times 3$
figurées par des points, puis un $\mathfrak{L}$ et une lettre biffée}

\[
  \begin{array}{l}
    (\mathrm{xiv}) \Rightarrow (\mathrm{i}) \Rightarrow (\mathrm{iii})
      \Rightarrow (\mathrm{v}) \\
    (\mathrm{xiv}) \searrow (\mathrm{x}?) \Rightarrow (\mathrm{xii})
  \end{array}
\]
\note{le numéro de la seconde ligne, noté ici $(\mathrm{x}?)$, est surchargé et
souligné, et ne se lit pas}

[Une seule racine $\alpha > 0$ telle que $\bar\alpha = 0$ ;
une racine $\beta > 0$ telle que $[X_{\alpha}, X_{\beta}] \neq 0$]

Alors $\mathfrak{L} + k\, \mathfrak{g}_{-\beta}$ est \uncertain{son}
\uncertain{propre} \uncertain{normalisateur}, \uncertain{satisfait} au th. de
densité, mais ne contient que des \ill{} algèbres de Cartan.\note{le tout
encadré ; « seule » et « satisfait » soulignés}

\end{document}
